% Generated draft: requires mathematical and source review.
\documentclass[11pt]{article}
\usepackage[margin=25mm]{geometry}
\usepackage{fontspec}
\setmainfont{Latin Modern Roman}
\usepackage{amsmath,amssymb,mathtools,unicode-math}
\setmathfont{Latin Modern Math}
\usepackage{enumitem,needspace,xurl,hyperref}
\setlength{\parindent}{0pt}
\setlength{\parskip}{5pt}
\setlength{\emergencystretch}{3em}
\newcommand{\N}{\mathbb{N}}
\newcommand{\Z}{\mathbb{Z}}
\newcommand{\Q}{\mathbb{Q}}
\newcommand{\R}{\mathbb{R}}
\newcommand{\C}{\mathbb{C}}
\newcommand{\F}{\mathbb{F}}
\newcommand{\A}{\mathbb{A}}
\newcommand{\PP}{\mathbb{P}}
\newcommand{\E}{\mathbb{E}}
\newcommand{\eps}{\varepsilon}
\newcommand{\dd}{\,\mathrm{d}}
\begin{document}
\section*{UnsolvedMath ID: 30003486}
\textbf{Source problem code:} catalogue entry
\subsection{Definitions and mathematical statement}
Let \(\mathcal M=\{\mathrm{CFN},\text{Jukes--Cantor},\mathrm{K2P},\mathrm{K3P}\}\) be the four model labels in the source. Write \(V_N^{(M)}\) for the variety denoted by \(V_N\) in the source when model \(M\in\mathcal M\) is used on a phylogenetic network \(N\).

\textbf{Requested statement.} For every \(M\in\mathcal M\) and every phylogenetic network \(N\) in the intended network class, determine defining polynomials for \(V_N^{(M)}\).

\textbf{Definitions needed to make the target precise.} The supplied record does not specify the network class, the model parameterizations, the ground field, the ambient coordinates, or whether the varieties are affine or projective. Consequently, it does not uniquely define \(V_N^{(M)}\). An affine formulation, conditional on supplying those choices, is as follows. Let \(k\) be the chosen algebraically closed field, let \(\Theta_{N,M}\) be the chosen parameter space, and let
\[
\phi_{N,M}:\Theta_{N,M}\longrightarrow k^{r_{N,M}}
\]
be the model parameterization in the chosen coordinates. Define
\[
V_N^{(M)}=\overline{\phi_{N,M}(\Theta_{N,M})}^{\mathrm{Zar}},\qquad
I(V_N^{(M)})=\{f\in k[x_1,\ldots,x_{r_{N,M}}]:f(p)=0\text{ for every }p\in V_N^{(M)}\}.
\]
Under these conventions, the requested equations would be a finite set \(F_{N,M}\subseteq k[x_1,\ldots,x_{r_{N,M}}]\) satisfying
\[
Z_k(F_{N,M})=V_N^{(M)},
\]
where \(Z_k(F)\) denotes the common zero set of the polynomials in \(F\). If the intended request is for generators of the full defining ideal, the requirement is instead
\[
\langle F_{N,M}\rangle=I(V_N^{(M)}).
\]
These requirements are distinct: equations defining the variety as a set need not generate its vanishing ideal. A projective interpretation would require a specified projective ambient space and homogeneous equations. The affine template above is conditional and does not resolve which conventions or which equation-generation target the source intends.
\subsection{Short English statement}
For each of the CFN, Jukes--Cantor, K2P, and K3P models, find polynomial equations describing the associated variety for every phylogenetic network. The supplied record leaves the network conventions, parameterizations, ambient space, and meaning of defining equations unspecified; these choices require clarification before there is a unique precise target.
\subsection{Sources}
\begin{itemize}
\item Local UnsolvedMath registry, supplied complete source record, permanent ID \(30003486\), title \emph{Defining Equations for Group-Based Phylogenetic Network Models}, source problem code \texttt{catalogue entry}, in the \emph{Oberwolfach Reports Open Problems} collection. No individual proposer is identified; the supplied proposed year is 2017.
\item Supplied report citation: \emph{Algebraic Statistics (2018)}, Oberwolfach Reports, DOI \texttt{10.4171/owr/2017/20}, pp.\ 1207--1279. \url{https://doi.org/10.4171/owr/2017/20}. The background attributes the problem collection to workshop open-problem discussions. The citation's 2018 label and the record's 2017 year are retained without resolving their relationship.
\item Supplied literature reference: \emph{Pfaffian structure for CFN networks}. \url{https://arxiv.org/abs/2312.07450}.
\item Supplied literature reference: \emph{Gross and Long on network ideals}. \url{https://doi.org/10.1137/17M1134238}.
\item Supplied literature reference: \emph{original OWR report}. \url{https://doi.org/10.4171/owr/2017/20}.
\item Literature assessment dated exactly \texttt{2026-08-21}: the supplied curation labels the problem \emph{partially solved} and reports that defining ideals are known for several small and structured network classes, including CFN level-one cases, while no general formula for arbitrary networks under all four models is known. This is supplied dated curation derived from the OWR collection, not an independently verified assessment. The supplied background explicitly calls for independent verification of time-sensitive status and cited claims.
\item The registry labels the statement \texttt{corrected\_verified}, explaining that mathematical content was preserved while notation, typography, and sentence structure were repaired. This supplied editorial label does not establish mathematical verification of the present draft.
\end{itemize}
\end{document}
